\defcap{sim01_timeseries.png}{腕 1 本の往復運動（下 → 上 → 下）における腕の角度と胴体の姿勢角（Sim01）。腕を戻すと姿勢も戻る。}
\defcap{sim01_sweep.png}{腕の質量比・長さ比と、180° の腕振りによる胴体の回転角（Sim01）。}
\defcap{sim02_cyclic.png}{左：曲率 $F(q)$ のマップと腕＋脚の長方形ループ。右：5 周分の胴体の姿勢角（Sim02）。}
\defcap{sim03_connection.png}{左：関節可動域内の曲率と最適円ループ。右：腕・脚・リアクションホイールの接続 $\Gamma$（Sim03）。}
\defcap{sim03_rest2rest.png}{90° の静止 → 静止旋回。左：リアクションホイール、右：AMBAC（Sim03）。}
\defcap{sim03_disturbance.png}{一定外乱トルクの下での姿勢保持と飽和（Sim03）。}
\defcap{sim04_human_cycle.png}{人間大モデルの順動力学（0.5 Hz の最適ループ）：姿勢、関節トルク、パワーと運動エネルギー（Sim04）。}
\defcap{sim04_scaling.png}{相似拡大による必要トルク・90° 旋回時間・角運動量容量の変化（Sim04）。}
\defcap{sim05a_learned_ambac.png}{学習した関節軌道と曲率マップ、姿勢の時系列、学習曲線（Sim05A）。}
\defcap{sim05b_pareto.png}{18 m 級の 90° 旋回における時間と推進剤のパレート曲線（左）と最適配分の例（右）（Sim05B）。}
\defcap{sim05b_robustness.png}{モデル誤差・外乱のあるプラントでの制御器の比較例（Sim05B）。}
\defcap{sim06_tipover.png}{荷物運び動作の ZMP（左上）、フルード則（右上）、軌道最適化（下）（Sim06）。}
\defcap{sim06_yaw_ratchet.png}{腕の水平振りによる台車のヨー（スティックスリップ）（Sim06）。}
\defcap{sim07_space.png}{宇宙での 3 次元 AMBAC：円錐運動による回転、回転軸、球面座標の特異点（Sim07）。}
\defcap{sim07_ground.png}{地上の 3 次元：定格荷重マップ、旋回中の ZMP、旋回中の腕の向き（Sim07）。}
\defcap{sim08_shoulder.png}{コッドマンのパラドックス（左 2 枚）と腕組みの姿勢・必要な内旋（右 2 枚）（Sim08）。}
\defcap{sim09_legs.png}{宇宙での脚の AMBAC 効果（左）、地上での戦車型と二本脚（中）、腰の戦略（右）（Sim09）。}
\defcap{sim10_steps.png}{段差の乗り越え：戦車型（左）、二本脚（中）、相似拡大（右）（Sim10）。}
